% CS F407 Artificial Intelligence: final project report template.
% Submission limit: 4 pages total.
% Do not change the font size, margins or spacing.
\documentclass[11pt,a4paper,twocolumn]{article}

\usepackage[margin=2cm,columnsep=0.8cm]{geometry}
\usepackage[T1]{fontenc}
\usepackage{lmodern}
\usepackage{microtype}
\usepackage{amsmath,amssymb}
\usepackage{graphicx}
\usepackage{booktabs}
\usepackage{tabularx}
\usepackage{enumitem}
\usepackage{titlesec}
\usepackage[hidelinks]{hyperref}

\titleformat{\section}{\normalsize\bfseries}{\thesection}{0.6em}{}
\titlespacing*{\section}{0pt}{1.2ex plus 0.3ex}{0.6ex}
\setlist{nosep,leftmargin=*}
\setlength{\parskip}{0.3ex}

\newcommand{\coursename}{CS F407 Artificial Intelligence, AY 2026--27 Semester I}
\newcommand{\teamno}{Neural Nexus}
\newcommand{\githuburl}{https://github.com/dhruvkshah75/groundtruth}

\begin{document}

\twocolumn[{%
\begin{center}
  {\large\bfseries GroundTruth: A Three-Layer Epistemic Architecture for Grounded Agents}\\[1ex]
  \teamno\\
  Dhruv Kenny Shah (2024A7PS0499G), Aaditya Lahori (2024A7PS0609G)\\
  Aayush Kushwaha (2024A7PS0524G), Rajat Surana (2024A1PS0291G)\\[0.5ex]
  {\small \coursename, BITS Pilani, K.K. Birla Goa Campus}\\[1.5ex]
  \begin{minipage}{0.85\textwidth}
    \small
    \centerline{\bfseries Abstract}\vspace{0.5ex}
    A robot answering from a static map can miss an obstacle. It can also
    confuse what a user expects with what its camera sees or what a
    maintenance log records. GroundTruth combines a provenance ledger and
    belief graph, a bounded language-model controller, and a simulated
    robot with viewpoint-limited sensors. The agent records why its beliefs
    change, keeps distinct perspectives visible, and reports uncertainty
    when it cannot verify the current state. We tested ten scripted
    scenarios in two offline modes: one replayed logged observations and
    the other reran the simulator and its sensors. All 30 query checks
    passed. The results show that deterministic evidence rules can make
    decisions about route status, object appearance, and
    location history traceable to the observations that support them.
  \end{minipage}
\end{center}
\vspace{1.5ex}
}]

\section{Introduction}
The course project, \emph{I, Agent}, asks how an agent should decide
what to believe when maps, sensors, and user expectations
disagree~\cite{assignment}. A map can be accurate when drawn and still
miss a cart that just blocked the route. A camera can describe how a
box looks under yellow light while a maintenance record describes
its painted color. The agent must know which questions each source
can answer.

GroundTruth takes a natural-language question in a selected scenario
or warehouse run. It returns an answer, an evidence status, and a
trace of supporting records. Users can ask about route status, when a
parcel was last seen, what changed its location belief, or why the
patrol stopped. We built a three-tier agent, a browser workspace for
inspecting the robot's evidence alongside the simulation, and a replay
suite that tests whether recorded observations reproduce the decisions.

\section{Background and related work}
The assignment draws on the Three-Tier Architecture (3T), which
separates deliberation, sequencing, and reactive skills in robot
control~\cite{bonasso1997}. GroundTruth adapts those boundaries to
declarative memory, procedural orchestration, and sensing.

The project plan also calls for ReAct, which interleaves model
reasoning with actions that retrieve external information~\cite{yao2023}.
GroundTruth uses the model to propose a typed intent and gives it the
evidence returned by the resulting tool plan. The project architecture
notes explain why structured queries and explicit rules govern belief
updates.

\section{Method}
\textbf{Tier boundaries.} Shared Pydantic contracts define the fact,
observation, intent, and result records. The procedural tier calls
memory and sensing through typed interfaces, evaluates evidence, and
requests justified revisions. The browser displays both the simulated
scene and the robot's evidence. A parcel behind a shelf can remain
visible to the viewer even when the robot cannot locate it.

\textbf{Declarative tier.} SQLite stores each claim's subject,
relationship, and value with source, confidence, observation time, and
spatial context. To revise a belief, the repository inserts a successor
claim, links its predecessor to it, and writes an audit event with the
evidence and rule in one transaction. Historical queries can still
retrieve the original claim and its source.

The rebuildable NetworkX multigraph uses fact IDs as edge keys to
preserve distinct active claims between the same nodes. The observation
ledger stores raw samples and indexes detections. A sample's run ID,
capture time, sequence, pose, and frame help distinguish what the robot
sensed and when from when the record reached storage.

\textbf{Procedural tier and runtime LLM.} Live mode uses Meta Llama 3.3
70B through Groq's function-calling API, configured as
\texttt{llama-3.3-70b-versatile}. The intent prompt asks the pretrained
model for one approved function call, with named sources separated
from entity mentions. For temporal queries it also selects a time
mode. Python validates the arguments, permits bounded repair,
resolves aliases, and checks sensor capabilities. An ambiguous match
prompts a clarification.

For a route question, the typed plan retrieves the stored claim and
front LiDAR evidence. The evaluator checks context and freshness before
updating the belief. For a color question it uses the detected object
ID to find the historical record and keeps user expectation, camera
appearance, and maintenance history separate. Python controls which
beliefs are written.

The explanation turn receives the verified results. Its prompt asks
for the conclusion, source attribution, and any uncertainty. Where
paraphrasing is allowed, a separate model call can check the proposed
wording against the evidence. Citation-sensitive queries require exact
wording. A rejected or unavailable explanation can yield a deterministic
answer or an error. Offline mode uses rules for intent and wording.

\textbf{Sensorimotor tier.} The mock environment computes LiDAR, camera,
light, and pose readings. The warehouse adds discrete time steps,
scheduled movement, and repeated sensing. Camera detections depend on
range, field of view, and occlusion. Front LiDAR measures the nearest
solid surface on its ray. Acquisition captures pose and readings
together. Playback speed changes presentation pacing while timestamps
follow simulation time.

The patrol subtracts distance traveled since the last LiDAR scan from
the measured clearance. It stops for a close obstacle and waits when a
scan is stale or unavailable. Its stop record retains the source sample
after the patrol resumes. A separate adapter validates recorded
telemetry and translates supported frames and timestamps for playback.

\textbf{Temporal evidence and questions.} A camera sighting inside one
station zone can establish a station belief. A new sighting at another
station can supersede it, while a repeated sighting in the same zone
refreshes verification. An aisle sighting can verify coordinates even
when the station belief lacks fresh support. A miss or camera outage
leaves earlier sightings available as history. A delayed older capture
cannot replace a newer applicable belief.

The controller answers current location, last seen, location at a
specified time, sighting history, and changes. For a cutoff $T$, it
reports the latest applicable capture at or before $T$ and gives the
actual observation time. Requests beyond the current run time return
uncertainty. Two sightings can establish a change while leaving its
path, exact timing, and cause undetermined.

\textbf{AI-assisted development.} Claude Code (Claude Opus 5.5)
and OpenAI GPT-6.2 Sol helped us write and revise the project code.
We reviewed the changes and checked behavior against project tests
and scenario outputs.

\section{Experimental setup}
\textbf{Problem instances.} Ten scripted scenarios specify initial
facts, sensor samples, and expected outcomes. They cover route
obstruction and clearance, visible relocation, movement out of view,
occlusion and reacquisition, ambiguous identity, color perspectives,
wrong-room or wrong-frame evidence, sensor outage and recovery, delayed
arrivals, and movement between captures. Some scenarios contribute
more than one query check.

\textbf{Replay modes.} Evidence-only replay ingests recorded
observations into a fresh ledger and answers using records available
up to each query checkpoint. A simulator rerun advances the world,
patrol, and sensors and checks regenerated evidence against the
recording. The wrong-context case includes a deliberately injected
sample. The full scene and presenter events stay outside the agent's
evidence inputs.

\textbf{Baseline and metrics.} The map-only comparison uses the
initial stored route status at the first cart obstruction. A replay
check passes when the status and all applicable assertions match.
Assertions can inspect answer content, location, active facts,
revision linkage, sample storage, or citation in the response trace.

\textbf{Execution environment.} We ran both offline modes on an Intel
Core i7-12650H with Python~3.12, SQLite, NetworkX, and Pydantic. They
require no GPU or Groq key. The React and TypeScript workspace has a
separate checkpoint test and demo runbook.

\section{Results}
All 30 query checks passed (Table~\ref{tab:results}). The checked-in
logs show expected and actual outcomes, applied policy, and supporting
evidence IDs.

\begin{table}[h]
\centering
\small
\begin{tabular}{lcc}
\toprule
Mode & Scenarios & Checks passed \\
\midrule
Evidence-only replay & 10 & 15/15 \\
Simulator rerun & 10 & 15/15 \\
\bottomrule
\end{tabular}
\caption{Offline replay checks. Higher pass counts are better.}
\label{tab:results}
\end{table}

\textbf{Groundedness and perspective.} In the replay route case,
the baseline map said clear. Front LiDAR detected the cart at
20~cm, and the agent stored a blocked claim. The audit showed
which observation superseded the old claim and why. A later clear
scan produced a new update. The separate assignment preset places
an obstacle at 12~cm. In the color case, both replay modes kept
red as the user's expectation, brown as the camera appearance, and
blue as the color recorded by maintenance bot \texttt{bot\_02}.
The simulator rerun also attributed the brown appearance to yellow
lighting. An ambiguous question about two boxes produced a request
for clarification.

\textbf{Location and time.} A visible relocation supported a
revision from the packing station to the inspection station.
The answer about the change cited the later camera sample and the
revision and left the route taken undetermined. When the parcel
moved out of view or a shelf occluded it, the agent retained the
last confirmed sighting and reported the current location as
uncertain. Reacquisition and camera recovery supplied new evidence.
The wrong-room and wrong-frame case left the warehouse's
stored location claim intact.

In the delayed-arrival case, a sample captured at 500~ms arrived
after the 2000~ms sighting. The inspection-station belief stayed
active and the answer cited the newer observation. In a different
case, a query about 500~ms fell between two captures. The agent
cited the 0~ms sighting and left the intermediate position
undetermined. The two cases check different ways time can matter:
the ordering of received records and the coverage of measurements.

\textbf{End-to-end demo.} At 6.5~s the parcel is behind a shelf.
The viewer can see it on the map while the last verified sighting
remains at Station A, captured at 6.0~s. By 12.5~s the camera
verifies aisle coordinates, and at 19.5~s the Station B detection
supports a station belief revision. The patrol stop at 28.0~s
links back to its LiDAR sample even after clearance and resumed
motion at the final 33.0~s checkpoint.
The route was first marked blocked at 25.5~s. The patrol continued
while the measured clearance still allowed motion, then withheld
movement when the obstacle reached the 50~cm threshold. This explains
why the route update and the stop have different evidence timestamps.

\section{Discussion and limitations}
The route baseline illustrates the benefit of checking a dynamic
claim against a current measurement at one checkpoint. The map
still provides useful geometry and initial context. The replay
counts measure agreement with prewritten expectations under scripted
conditions. Evidence-only replay checks the decisions without the
simulator, while the rerun also checks whether the runtime and
sensors reproduce the observations. A benchmark spanning more
layouts and unscripted motion would test how broadly the results hold.

The mock camera uses known object IDs and simplified geometry.
The color case applies a fixed lighting rule and a seeded
maintenance record. It tests whether the agent attributes the three
claims correctly under a known conflict. An actual camera can miss or
misidentify objects for additional reasons. A verified sample can
also become stale as soon as an object moves, and a missed
detection tells the agent little about where the object went.

The same issue applies to clearance. A LiDAR scan establishes the
space observed in the tested direction at capture time. When the robot
turns, the previous scan cannot verify forward clearance on the new
heading. The simulator can exercise that rule with known geometry.
Testing it with recorded or live measurements would also require
checking that sensor frames, capture times, and robot poses are aligned.

The reported evaluation runs offline. A live model can misclassify
unfamiliar wording, and the model-based check of a paraphrase can
miss an unsupported claim. The typed plans and evidence rules still
constrain which operations run and which beliefs can be written.
Next steps include held-out questions phrased by different users,
real recordings with noise and delayed samples, and testing the
calibrated telemetry pipeline before using it on a physical robot.
For the live model, those tests could score intent selection and
answer fidelity separately. An incorrectly selected intent can
retrieve the wrong kind of evidence, even when the resulting answer
faithfully describes what was retrieved.

\section{Conclusion}
GroundTruth routes questions through a bounded language interface
and ties its answers and belief revisions to identifiable evidence.
Across ten scripted scenarios, both replay modes reproduced the
expected conflict resolution, perspective tracking, and uncertainty
behavior. The warehouse demo makes those decisions inspectable by
linking temporal answers, location revisions, and patrol stops to the
records available to the robot.

\section*{Individual member contributions}
{\small
\begin{tabularx}{\linewidth}{@{}l X@{}}
\toprule
Name & Contribution \\
\midrule
Dhruv Shah & Designed shared contracts and intent planning, including
entity and capability checks. Built the warehouse engine and sensors
and integrated the session runtime, HTTP API, and React workspace.
Connected the belief timeline to the demo, verified its checkpoints,
and finalized the replay results. \\
Aaditya Lahori & Built the NetworkX active-belief projection, its lazy
cache, and consistency checks against SQLite. Implemented temporal
ingestion that turns comparable camera and LiDAR observations into
auditable belief revisions. Added safeguards for late captures,
run isolation, and retries. \\
Aayush Kushwaha & Implemented the SQLite fact and audit ledger,
atomic revisions, and graph cache invalidation. Built the core
epistemic evaluator and Groq function-calling integration. Implemented
the sensor-backed patrol controller, ten-scenario export and replay
suite, and the recorded telemetry adapter. \\
Rajat Surana & Built the original mock environment and the typed
plan executor and composition service. Implemented the raw
observation ledger and its run-scoped queries, added temporal location
and change questions, and refined off-station parcel verification and
the LiDAR evidence attached to patrol stops. \\
\bottomrule
\end{tabularx}}

\smallskip
\noindent Contributions reflect component work and integration recorded in
the \href{https://github.com/dhruvkshah75/groundtruth/commits/50127947154341efce02b63e7e859074764ea1fe/}{merged commit history}.
All members have read this report and agree with the statement above.

\section*{Generative AI disclosure}
We used Anthropic Claude Code (Claude Opus 5.5) and OpenAI GPT-6.2 Sol
to help write the project code. We checked the AI-assisted work and
take responsibility for the submitted code and report.

\section*{Code availability}
\noindent Code and data:\\
\url{\githuburl}\\
The public repository contains the agent, scripted scenarios, tests,
and evaluation artifacts. To reproduce Table~\ref{tab:results}, use
commit \texttt{5012794}, run \verb|uv sync --locked| from the repository
root, then run the command below. The README's documentation index
links to the \href{https://github.com/dhruvkshah75/groundtruth/blob/50127947154341efce02b63e7e859074764ea1fe/docs/phase-two/Replay_and_Evaluation.md}{replay guide}.

\begin{verbatim}
uv run --locked python -m \
  src.evaluation.cli --mode both \
  --output-dir /tmp/groundtruth-eval
\end{verbatim}

The command writes the per-query Markdown report and JSON summary.
The published copies are under \path{artifacts/evaluation/}. The
demo runbook provides the offline browser setup and exact checkpoints
for inspecting the warehouse timeline.

\begin{thebibliography}{9}
\small
\bibitem{assignment}
CS F407 Artificial Intelligence.
\newblock \emph{I, Agent (A Three-Layer Epistemic Architecture for Grounded Agents)}.
\newblock Project problem statement, BITS Pilani, K.K. Birla Goa Campus, 2026.
\bibitem{bonasso1997}
R.~P. Bonasso, R.~J. Firby, E.~Gat, D.~Kortenkamp, D.~P. Miller, and M.~G. Slack.
\newblock \href{https://doi.org/10.1080/095281397147103}{Experiences with an architecture for intelligent, reactive agents}.
\newblock \emph{Journal of Experimental \& Theoretical Artificial Intelligence}, 9(2--3):237--256, 1997.
\bibitem{yao2023}
S.~Yao, J.~Zhao, D.~Yu, N.~Du, I.~Shafran, K.~Narasimhan, and Y.~Cao.
\newblock \href{https://arxiv.org/abs/2210.03629}{ReAct: Synergizing reasoning and acting in language models}.
\newblock \emph{ICLR}, 2023.
\end{thebibliography}

\end{document}
